\section{Effect of Disturbance Geometry on Robust Steering}
\label{app:disturbance-geometry}

We examine how the disturbance geometry used for robust synthesis affects
steering quality. Using Qwen2.5-0.5B-Instruct on HarmBench, we hold the
eight-dimensional state representation, fitted dynamics, steering target,
and costs fixed and change only the disturbance covariance used for
synthesis. At each layer, we compare the empirical covariance $C_k$ with a
randomly rotated covariance $UC_kU^\top$ and an isotropic covariance
$\operatorname{tr}(C_k)I/8$, where $U$ is a fixed orthogonal matrix.
All three preserve total disturbance variance; rotation additionally
preserves the covariance spectrum while changing its orientation.
Each controller is synthesized using the same $H_\infty$ procedure.
Performance is pooled over 200 direct prompts and 200 jailbreak prompts,
with the same generation and scoring pipeline used in the main experiments.

\begingroup
\renewcommand{\thefigure}{R2}
\begin{figure}[H]
    \centering
    \includegraphics[width=0.80\linewidth]
        {Final_Figures/geometry_final_for_paper.pdf}
    \caption{
    \textbf{Disturbance geometry and robust steering on
    Qwen2.5-0.5B-Instruct.}
    \textbf{a}: AXBench overall for empirical, randomly rotated,
    and isotropic disturbance geometries.
    \textbf{b}: Normalized Frobenius alignment with held-out OOD
    residual covariances. Dots identify layers; boxes summarize
    their distribution. Statistical significance was assessed using
    paired bootstrap tests with Bonferroni correction
    ($^{***}p<0.001$).
    }
    \label{fig:disturbance-geometry}
\end{figure}
\endgroup

Empirical geometry provides the highest steering quality among the three
robust controllers: AXBench overall is $0.440$ for empirical geometry and
$0.352$ after rotation (Figure~\ref{fig:disturbance-geometry}a).
This joint score combines concept relevance, instruction relevance, and
fluency. Rotation therefore reduces steering quality despite preserving
the disturbance magnitude and covariance spectrum.
We measure layer-wise covariance alignment using the normalized Frobenius
inner product
$\operatorname{Align}_k=\langle\widetilde C_k,\widehat C_k^{\mathrm{OOD}}\rangle_F/(\|\widetilde C_k\|_F\|\widehat C_k^{\mathrm{OOD}}\|_F)$,
where $\widetilde C_k$ is the synthesis covariance derived from the
calibration estimate $C_k=\widehat{\Sigma}_k$ and
$\widehat C_k^{\mathrm{OOD}}$ is the held-out OOD residual covariance.
An alignment of one indicates proportional covariance matrices.
Empirical geometry also has the highest alignment with held-out OOD
residual covariances, with mean alignment decreasing from $0.818$ to
$0.416$ under rotation (Figure~\ref{fig:disturbance-geometry}b).
These results show that the empirical orientation of model mismatch
contributes to the quality of robust steering.

Replacing the covariance with an isotropic model produces a more severe
degradation. Despite preserving total disturbance variance, isotropic
synthesis yields largely empty or corrupted responses and an AXBench
overall score of $0.005$. Retaining disturbance magnitude alone is
therefore insufficient to preserve useful generation.
Together, these comparisons show that steering quality depends on how
disturbance variance is distributed and oriented within the state space.
Empirical disturbance geometry preserves the measured
directional structure of model mismatch and yields higher steering
quality than either alternative under the same robust synthesis procedure.
